% =====================================================================
\subsection{ETSI GS QKD 015: the SDN controller and the SD-QKD network}
\label{sec:qkd015}
% =====================================================================

GS~QKD~015 (V1.1.1 in 2021, V2.1.1 in April 2022)~\cite{qkd015} defines
the control interface between an SDN controller and the nodes of a QKD
network, i.e.\ the \emph{southbound interface} of the QKD control plane.
The second edition, considered here, was aligned with GS~QKD~018: the
GS~QKD~018 YANG modules import those of GS~QKD~015 directly. A good
understanding of 015 is therefore a prerequisite for understanding 018.

The section proceeds in four steps: why software control is needed; what
the components of an SD-QKD network are, one by one; what a YANG model
looks like in general; and how 015 uses YANG to describe the node, with the
resulting workflows and security considerations.

% ---------------------------------------------------------------------
\subsubsection{Why software-defined control}

The specification starts from an observation: integrating QKD into a
telecommunications network is today a manual operation. Each new QKD
system is installed by hand, usually over a dedicated dark fibre between
two sites, forming a fixed link. This approach has three limits. It does
not scale, because every change requires human intervention. It wastes
resources, because the key capacity of a link cannot be moved to where it
is needed. And it does not talk to the optical network, which is
increasingly disaggregated and software-controlled.

The proposed architecture lets a controller orchestrate QKD resources
centrally: optimise the key allocation per link according to demand,
automatically create direct or multi-hop links, and coordinate quantum and
classical channels even when they share the same fibre. In particular, the
abstraction defined by the specification aims to:
\begin{itemize}
  \item manage QKD resources centrally according to demands, capabilities
        and network state;
  \item aggregate the QKD systems of the same perimeter under a single key
        management, to better detect demands;
  \item hide the complexity of the low-level parameters of each technology,
        delegating them to the node;
  \item optimise the distribution of links according to application QoS and
        generation statistics;
  \item coordinate quantum and classical channels, whether or not they
        share the same physical medium.
\end{itemize}
These are, in technical form, the needs that emerged from use cases UC2,
UC3 and UC4 of GS~QKD~002.

% ---------------------------------------------------------------------
\subsubsection{The components of an SD-QKD network}

It is convenient to build the network bottom-up, adding one abstraction at
a time. Fig.~\ref{fig:sdqkd} shows the result.

\paragraph{1. QKD module, system and link.}
The basic physical unit is the \emph{QKD module}: the set of hardware and
software that runs part of a QKD protocol. Two connected modules form a
\emph{QKD system}, and the connection between them (quantum channel plus
classical channel) is the \emph{QKD link}. So far we are in the quantum
layer of Fig.~\ref{fig:itu}: modules are vendor- and technology-specific
(discrete or continuous variables, encoding, wavelength) and do not
interoperate.

\paragraph{2. QKD interface.}
The controller does not need to know how a module works internally, only
what it offers to the network. GS~QKD~015 therefore introduces an
abstraction.

\begin{definition}[QKD interface]
A \emph{QKD interface} (\texttt{qkdi}) is the abstract representation of a
QKD module within a node: it exposes only the attributes relevant to the
network, such as technology, model, wavelength range, maximum supported
attenuation and attachment point to an optical switch.
\end{definition}

The analogy is the ports of a switch: a node with three QKD modules appears
to the controller as a device with three interfaces. Since modules from
different vendors do not interoperate, the model requires the technology
and model of each interface to be declared: the controller will pair only
compatible interfaces.

\paragraph{3. SD-QKD node.}
Interfaces located in the same secure site are grouped together.

\begin{definition}[SD-QKD node]
An \emph{SD-QKD node} is the aggregation of one or more QKD modules located
in a secure site and managed by a single \emph{SDN agent}, which exposes
them to the controller as a single network element with several interfaces.
\end{definition}

The specification sets a series of requirements that give the node its
shape. An SD-QKD node \emph{shall}: comprise at least one QKD module,
exposed as an interface; be located in a secure site (a site may host
several nodes); contain a \emph{key management system} that aggregates the
material produced by all key associations; offer applications a
\emph{single access point} to keys through an API; and connect to at least
one SDN controller via standard protocols for discovery, control and
telemetry. It \emph{should} also expose to the controller information on
applications, interfaces, links and classical-channel requirements (for
example acceptable attenuation).

Inside the node four logical components can thus be identified:
\begin{itemize}
  \item the \textbf{QKD modules}, which generate key with remote nodes;
  \item the \textbf{key manager} (KMS), which receives keys from the
        modules, stores them in one or more logical \emph{key stores},
        relays them when needed and registers applications with their QoS;
  \item the \textbf{application interface}, i.e.\ the key delivery API
        (GS~QKD~004 or~014), reachable at a known IP address and port;
  \item the \textbf{SDN agent}, which translates the state of all of the
        above into a YANG model, exposes it to the controller, receives
        configurations and sends notifications.
\end{itemize}

% Figure: SD-QKD network and internal structure of a node (adapted from GS QKD 015, Fig. 1)
\begin{figure}[t]
\centering
\begin{tikzpicture}[
  comp/.style={draw,rounded corners,align=center,font=\scriptsize,
               minimum height=6mm},
  small/.style={draw,circle,fill=green!12,minimum size=9mm,font=\scriptsize,
                align=center}]
  % controller
  \node[box,fill=blue!8,minimum width=4cm] (ctl) at (-0.6,4.3) {SDN controller (SD-QNC)};
  % detailed node
  \node[comp,fill=yellow!20,minimum width=1.4cm] (m1) at (-2.2,0) {QKD\\module 1};
  \node[comp,fill=yellow!20,minimum width=1.4cm] (m2) at (-0.6,0) {QKD\\module 2};
  \node[comp,fill=yellow!20,minimum width=1.4cm] (m3) at (1.0,0) {QKD\\module 3};
  \node[comp,fill=green!12,minimum width=4.6cm] (kms) at (-0.6,1.0) {key manager (KMS, key store)};
  \node[comp,fill=blue!6,minimum width=1.9cm] (agent) at (-1.95,2.0) {SDN agent};
  \node[comp,fill=gray!10,minimum width=2.4cm] (app) at (0.75,2.0) {applications\\(QKD 004/014)};
  \begin{scope}[on background layer]
    \node[draw,dashed,rounded corners,fill=gray!3,inner sep=6pt,
          fit=(m1)(m3)(kms)(agent)(app),
          label={[font=\scriptsize\itshape]below:SD-QKD node (secure perimeter)}] (node) {};
  \end{scope}
  % remote nodes
  \node[small] (r1) at (4.6,2.6) {node\\B};
  \node[small] (r2) at (4.4,-0.2) {node\\C};
  \node[small] (r3) at (-5.2,0.6) {node\\D};
  % links
  \draw[thick] (m3.east) -- node[lbl,above]{PHYS} (r2);
  \draw[thick] (m1.west) -- (r3);
  \draw[thick] (r1) -- (r2);
  \draw[thick,dotted,red!70!black] (kms.east) -- node[pos=0.55,font=\scriptsize,fill=white,inner sep=1pt]{VIRT via C} (r1);
  \draw[<->] (m1) -- (m1|-kms.south);
  \draw[<->] (m2) -- (m2|-kms.south);
  \draw[<->] (m3) -- (m3|-kms.south);
  \draw[<->] (agent.south) -- (agent.south|-kms.north);
  \draw[<->] (app.south) -- (app.south|-kms.north);
  \draw[<->,dashed,blue!60!black] (agent.north) -- node[pos=0.5,left,font=\scriptsize]{GS QKD 015} (ctl);
  \draw[<->,dashed,blue!60!black] (r1) -- (ctl);
  \draw[<->,dashed,blue!60!black] (r3) -- (ctl);
\end{tikzpicture}
\caption{An SD-QKD network: nodes are connected by physical (solid) and
virtual (dotted) links, and each is controlled by the controller through its
SDN agent (dashed). The node in the foreground shows the internal
components: the QKD modules, exposed as interfaces, the key manager, the
application interface and the agent. Adapted from~\cite[Fig.~1]{qkd015}.}
\label{fig:sdqkd}
\end{figure}


\paragraph{4. Key association link.}
Between two nodes the interesting quantity is not the quantum channel
itself but the fact that the two nodes share key. GS~QKD~015 models exactly
this.

\begin{definition}[Key association link]
A \emph{key association link} (\texttt{qkd\_link}) is a logical key
association between two remote SD-QKD nodes. It is of type \texttt{PHYS}
(direct) if a quantum channel connects the two nodes, of type \texttt{VIRT}
(virtual) if the key is obtained by key relay through intermediate trusted
nodes.
\end{definition}

A physical link corresponds to a QKD link in ITU-T terminology; a virtual
link corresponds to a chain of KM links over several QKD links. Two details
make the model consistent. First, a virtual link \emph{consumes} key from
the physical links it rests on, and is therefore registered as an
\emph{internal application} with a reserved bandwidth
(\texttt{virt\_bandwidth}). Second, every intermediate node of a virtual
link knows the previous and the next node (\texttt{virt\_prev\_hop},
\texttt{virt\_next\_hop}), so the relay path is configured hop by hop by
the controller. For a node to act as a relay it must declare the
\texttt{key\_relay\_mode\_enable} capability.

\paragraph{5. Applications.}
From the node's point of view, an application is any entity that requests
keys from the key manager.

\begin{definition}[QKD application]
A \emph{QKD application} (\texttt{qkd\_app}) is a running instance that
consumes keys from an SD-QKD node. It is \emph{external}
(\texttt{CLIENT}) if it is a user, an HSM, an encryptor, a virtual network
function or a security protocol; it is \emph{internal}
(\texttt{INTERNAL}) if it uses keys for the operation of the network
itself, such as authentication or relay for a virtual link.
\end{definition}

Each application declares a quality of service: maximum and minimum key
bandwidth (\texttt{max\_bandwidth}, \texttt{min\_bandwidth}, in bit/s),
maximum delivery \texttt{jitter} in ms, lifetime \texttt{ttl} of an unused
key in the key store, and whether paths or keys may be shared among several
clients. This is the technical answer to the differentiated-QoS need of
UC2.

\paragraph{6. SDN controller.}
At the top sits the controller (in 015 also \emph{SD-QNC},
\emph{Software-Defined QKD Network Controller}), which has the global view
of the domain and plays three roles. It acts as the \emph{central registry}
of applications, so that an application only needs to know its own key
manager and not the network topology. It plans and creates \emph{links},
physical or virtual, according to demand and statistics. And it
coordinates its domain with the optical network controller
(\emph{SD-ONC}) when a new optical path is needed for a quantum channel.

\paragraph{What does not reach the controller.}
The information that crosses the control plane is deliberately minimal:
identifiers, states, rates, physical parameters. Key material \emph{never
leaves the node perimeter}. For this reason, the specification observes,
introducing SDN into a trusted-node network adds no new risks beyond those
already known, such as denial of service.

% ---------------------------------------------------------------------
\subsubsection{Anatomy of a YANG model}
\label{sec:yang-anatomia}

Both 015 and 018 are written in YANG~\cite{rfc7950}, a language born in the
IETF to describe the configuration and state data of network devices and
the operations that can be invoked on them. A YANG model is not executable
code: it is a \emph{schema}, like an XSD for XML or a JSON schema, which
states which data exist, of which type, how they are nested and who may
modify them. The NETCONF and RESTCONF protocols then work on top of that
schema. We go through the main constructs using examples taken from the
015 files~\cite{qkd015yang}.

\paragraph{The module and its header.}
The basic unit is the \emph{module}. Its header declares the language
version, a globally unique \texttt{namespace}, a short \texttt{prefix} used
by other modules to refer to it, the imported modules and the revisions.
\begin{lstlisting}[language=yang]
module etsi-qkd-sdn-node {
  yang-version "1";
  namespace "urn:etsi:qkd:yang:etsi-qkd-node";
  prefix "etsi-qkdn";
  import ietf-yang-types   { prefix "yang"; }
  import ietf-inet-types   { prefix "inet"; }
  import etsi-qkd-node-types { prefix "etsi-qkdn-types"; }
  organization "ETSI ISG QKD";
  revision "2022-01-30" { reference "ETSI GS QKD 015 V2.1.1"; }
  ...
}
\end{lstlisting}
GS~QKD~015 is split into two modules. The first one,
\texttt{etsi-qkd-node-types}, contains only types; the second one,
\texttt{etsi-qkd-sdn-node}, contains the actual model. This is
a common choice: separating the types lets other modules, such as those of
018, reuse them without importing everything else.

\paragraph{Types: built-in, typedef and identity.}
YANG offers built-in types (\texttt{boolean}, \texttt{string},
\texttt{uint32}, \texttt{decimal64}, \texttt{enumeration}\dots) and common
types defined by the IETF, such as \texttt{yang:uuid} or
\texttt{yang:date-and-time}. \texttt{typedef} defines derived types with
their own name. For extensible symbolic values, \texttt{identity} is used:
a base identity is declared, then identities that derive from it; a field
of type \texttt{identityref} accepts any identity derived from the base.
This is how 015 defines link types:
\begin{lstlisting}[language=yang]
identity QKD-LINK-TYPES { description "QKD key association link types."; }
identity VIRT { base QKD-LINK-TYPES; description "Virtual Link."; }
identity PHYS { base QKD-LINK-TYPES; description "Physical Link."; }

typedef qkd-link-types {
  type identityref { base QKD-LINK-TYPES; }
}
\end{lstlisting}
The advantage over an \texttt{enumeration} is that another module can add,
say, a new link type without modifying the original. Technologies
(\texttt{CV-QKD}, \texttt{DV-QKD}\dots), node states (\texttt{NEW},
\texttt{OPERATING}, \texttt{DOWN}, \texttt{FAILURE}, \texttt{OUT}), roles
(\texttt{TRANSMITTER}, \texttt{RECEIVER}, \texttt{TRANSCEIVER}) and
application types (\texttt{CLIENT}, \texttt{INTERNAL}) are defined in the
same way.

\paragraph{Data nodes: container, list, leaf, leaf-list.}
Data are organised as a tree with four constructs:
\begin{itemize}
  \item \texttt{leaf}: a single typed value (e.g.\ \texttt{qkdn\_id});
  \item \texttt{leaf-list}: a sequence of values of the same type (e.g.\
        \texttt{virt\_next\_hop}, which allows several successors);
  \item \texttt{container}: a node grouping other nodes, with no value of
        its own (e.g.\ \texttt{qkdn\_capabilities});
  \item \texttt{list}: a collection of structured entries, each identified
        by a \texttt{key} (e.g.\ the \texttt{qkd\_link} list keyed by
        \texttt{qkdl\_id}).
\end{itemize}

\paragraph{Configuration and state.}
Every node is either \emph{configuration} (\texttt{config true}, the
default: the client may write it) or \emph{operational state}
(\texttt{config false}: produced by the device, read-only for the client).
In 015, for example, \texttt{qkdl\_enable} is configurable, while
\texttt{qkdn\_status}, \texttt{skr} and \texttt{eskr} are state. This
distinction mirrors the split between what the controller \emph{decides}
and what the node \emph{measures}.

\paragraph{Reuse and extension: grouping, uses, augment.}
A \texttt{grouping} is a reusable block of nodes, copied wherever
\texttt{uses} appears. The 015 model is built almost entirely this way:
\texttt{app\_basic}, \texttt{app\_details}, \texttt{app\_qos},
\texttt{qkd\_links} and so on are groupings later assembled into the
\texttt{qkd\_node} container. \texttt{augment} does the opposite: it adds
nodes to a tree defined \emph{in another module}, without touching it. It
is the tool that 018 will use to enrich its topology with inventory data.

\paragraph{References and constraints.}
A \texttt{leafref} is a field whose value must match an existing value
elsewhere in the tree, like a foreign key in a database. \texttt{when} and
\texttt{must} statements add conditions, for instance that a field exists
only if the link is of a certain type. This is exactly how 015 states that
the \texttt{virt\_*} fields apply only to virtual links and the
\texttt{phys\_*} fields only to physical ones.

\paragraph{Operations and notifications: rpc and notification.}
Besides data, a module can define operations invoked by the client
(\texttt{rpc}, with \texttt{input} and \texttt{output}) and messages the
server sends spontaneously (\texttt{notification}). 015 mostly defines
notifications; 018, as we will see, also defines RPCs.

\paragraph{Reading a tree.}
The \texttt{pyang} tool produces a tree view of a module, which we use
below. The notation is compact: \texttt{rw} marks a configuration node,
\texttt{ro} a state node, \texttt{?} an optional node, \texttt{*} a list or
leaf-list, a list's key appears in square brackets and \texttt{-{}-{}-n}
marks a notification.

\paragraph{From data to message.}
The same schema produces both XML and JSON. In JSON~\cite{rfc7951} the top
level is qualified with the module name, and \texttt{identityref} values
carry the module that defines the identity. A physical link exposed by a
node would look like this:
\begin{lstlisting}
{ "etsi-qkd-sdn-node:qkd_node": {
    "qkdn_id": "0f3c...-a1",
    "qkd_links": { "qkd_link": [ {
        "qkdl_id": "7b2e...-c4",
        "qkdl_type": "etsi-qkd-node-types:PHYS",
        "qkdl_local":  { "qkdn_id": "0f3c...-a1", "qkdi_id": 1 },
        "qkdl_remote": { "qkdn_id": "9d41...-e7", "qkdi_id": 2 },
        "phys_wavelength": 1550,
        "qkdl_performance": { "skr": 20000, "eskr": 12000,
                              "expected_consumption": 8000 } } ] } } }
\end{lstlisting}
(Values are illustrative.)

% ---------------------------------------------------------------------
\subsubsection{The SD-QKD node model}

With these tools the 015 model is easy to read. The root container
\texttt{qkd\_node} has four sections, one for each component described
above, plus the notifications:
\begin{lstlisting}
module: etsi-qkd-sdn-node
  +--rw qkd_node
     +--rw qkdn_id              yang:uuid
     +--ro qkdn_status?         etsi-qkdn-types:qkdn-status-types
     +--rw qkdn_version?        string
     +--rw qkdn_location_id?    string
     +--rw qkdn_capabilities
     |  +--rw link_stats_support?          boolean
     |  +--rw application_stats_support?   boolean
     |  +--rw key_relay_mode_enable?       boolean
     +--rw qkd_applications
     |  +--rw qkd_app* [app_id]
     |     +--rw app_id, app_type, server_app_id, client_app_id*, ...
     |     +--rw app_qos            (max/min_bandwidth, jitter, ttl, ...)
     |     +--rw backing_qkdl_id*, local_qkdn_id, remote_qkdn_id
     +--rw qkd_interfaces
     |  +--rw qkd_interface* [qkdi_id]
     |     +--rw qkdi_id, qkdi_model, qkdi_type
     |     +--rw qkdi_att_point     (device, port)
     |     +--rw qkdi_capabilities  (role, wavelength_range, max_absorption)
     |     +--ro qkdi_status
     +--rw qkd_links
        +--rw qkd_link* [qkdl_id]
           +--rw qkdl_id, qkdl_status, qkdl_enable, qkdl_type
           +--rw qkdl_local / qkdl_remote   (qkdn_id, qkdi_id)
           +--rw qkdl_applications*
           +--rw virt_prev_hop, virt_next_hop*, virt_bandwidth   (VIRT)
           +--rw phys_channel_att, phys_wavelength, phys_qkd_role (PHYS)
           +--rw qkdl_performance   (expected_consumption, skr, eskr,
                                     phys_perf* [perf_type]: QBER, SNR)
\end{lstlisting}

\paragraph{Node identity.}
\texttt{qkdn\_id} is a UUID unique in the network; \texttt{qkdn\_status},
\texttt{qkdn\_version} and \texttt{qkdn\_location\_id} serve inventory and
diagnostics.

\paragraph{Capabilities.}
The three capabilities tell the controller what the node is willing to do:
expose link statistics (SKR, consumption, status, QBER), expose
application statistics, and act as a relay. A node with
\texttt{key\_relay\_mode\_enable = false} will never be chosen as an
intermediate node of a virtual link.

\paragraph{Applications.}
Each \texttt{qkd\_app} binds an application to the key service:
\texttt{server\_app\_id} identifies the entity that initiated the request,
\texttt{client\_app\_id} the entities allowed to receive the same key,
\texttt{backing\_qkdl\_id} the link that supplies the keys,
\texttt{local\_qkdn\_id} and \texttt{remote\_qkdn\_id} the two nodes
involved. Until the remote node is known, the local node cannot serve keys.
Statistics (\texttt{app\_statistics}) report the bits consumed per time
interval.

\paragraph{Interfaces.}
Each \texttt{qkd\_interface} describes a module: model, technology, status,
attachment point to an optical device (\texttt{qkdi\_att\_point}) and
physical capabilities. This is the information the controller needs to
pair two compatible interfaces and request an optical path between their
ports.

\paragraph{Links and performance.}
Each \texttt{qkd\_link} reports its endpoints, type, status, the
applications it serves and the type-specific parameters. The
\texttt{qkdl\_performance} block holds the three quantities that drive all
planning:
\begin{itemize}
  \item \texttt{skr}: secret key generation rate, in bit/s;
  \item \texttt{expected\_consumption}: sum of the bandwidths of all
        applications on the link;
  \item \texttt{eskr}: \emph{effective} rate, i.e.\ what is available
        after internal consumption.
\end{itemize}
In compact form, if $\mathcal{A}_\ell$ is the set of applications on link
$\ell$, each with bandwidth $b_a$,
\begin{equation}
  \text{expected\_consumption}_\ell = \sum_{a\in\mathcal{A}_\ell} b_a ,
  \qquad
  \mathrm{ESKR}_\ell \approx \mathrm{SKR}_\ell - \text{expected\_consumption}_\ell .
  \label{eq:eskr015}
\end{equation}
The second relation is an intuitive reading of the fields, not a formula of
the standard. Physical links also carry metrics such as QBER and SNR.

% ---------------------------------------------------------------------
\subsubsection{Notifications}

Polling the nodes continuously would be inefficient. GS~QKD~015 therefore
defines YANG notifications from node to controller, grouped by object:
\begin{itemize}
  \item \textbf{applications}: \texttt{sdqkdn\_application\_new} (a new
        application requests keys), \texttt{sdqkdn\_application\_qos\_update},
        \texttt{sdqkdn\_application\_disconnected};
  \item \textbf{interfaces}: \texttt{sdqkdn\_interface\_new},
        \texttt{sdqkdn\_interface\_down} (with a \emph{reason} field),
        \texttt{sdqkdn\_interface\_out};
  \item \textbf{links}: \texttt{sdqkdn\_link\_down},
        \texttt{sdqkdn\_link\_perf\_update} (significant performance change),
        \texttt{sdqkdn\_link\_overloaded} (the link cannot cope with demand).
\end{itemize}
In addition there is a generic structure for events and alarms, designed to
accommodate new event types without changing the model. Moreover,
YANG-based protocols allow subscribing to any datum in the datastore, so
that every change can be reported in real time.

% ---------------------------------------------------------------------
\subsubsection{Workflows}

Clause 5 describes three fundamental scenarios that show the components at
work.

\paragraph{Application registration.}
This is the most instructive case, because it explains why the controller
acts as a central registry (Fig.~\ref{fig:seq015}). An application knows
only the address of its own key manager; it does not know the node of the
application it wants to talk to.
\begin{enumerate}
  \item The application in security domain $A$ connects to its node
        through the key delivery API, sending its own ID, the ID of the
        peer application and the requested QoS.
  \item Node $A$ does not know where the peer application is and forwards
        the request to the controller, adding its own ID.
  \item If the peer application is not yet registered, the controller
        replies with a plain acknowledgement.
  \item In domain $B$ the peer application connects to its node, which in
        turn informs the controller.
  \item The controller matches the pair, generates a unique service
        identifier (\texttt{app\_id}) and sends both nodes all the
        information they need: application and node IDs, \texttt{app\_id}
        and QoS.
  \item The nodes pass the \texttt{app\_id} to the applications and reserve
        keys for them; from then on the applications request keys using
        that identifier.
\end{enumerate}

% Figure: application registration (adapted from GS QKD 015, Fig. 2)
\begin{figure}[t]
\centering
\begin{tikzpicture}[font=\scriptsize]
  \def\xA{0}\def\xB{3.0}\def\xC{6.6}\def\xD{10.2}\def\xE{13.2}
  \node[actor,minimum width=2cm] at (\xA,0) {App A};
  \node[actor,minimum width=2cm,fill=green!8] at (\xB,0) {Node A};
  \node[actor,minimum width=2cm,fill=blue!8] at (\xC,0) {Controller};
  \node[actor,minimum width=2cm,fill=green!8] at (\xD,0) {Node B};
  \node[actor,minimum width=2cm] at (\xE,0) {App B};
  \foreach \x in {\xA,\xB,\xC,\xD,\xE} \draw[gray] (\x,-0.4) -- (\x,-6.6);
  \draw[msg] (\xA,-0.9) -- node[above]{ID, ID peer, QoS} (\xB,-0.9);
  \draw[msg] (\xB,-1.5) -- node[above]{\texttt{sdqkdn\_application\_new}} (\xC,-1.5);
  \draw[ret] (\xC,-2.1) -- node[above]{ack (peer unknown)} (\xB,-2.1);
  \draw[msg] (\xE,-2.9) -- node[above]{ID, ID peer, QoS} (\xD,-2.9);
  \draw[msg] (\xD,-3.5) -- node[above]{\texttt{sdqkdn\_application\_new}} (\xC,-3.5);
  \draw[msg] (\xC,-4.1) -- ++(0.35,0) -- ++(0,-0.35) -- (\xC,-4.45);
  \node[right] at (\xC+0.4,-4.27) {match and create \texttt{app\_id}};
  \draw[msg] (\xC,-5.1) -- node[above]{config.\ (\texttt{app\_id}, QoS)} (\xB,-5.1);
  \draw[msg] (\xC,-5.1) -- node[above]{config.\ (\texttt{app\_id}, QoS)} (\xD,-5.1);
  \draw[ret] (\xB,-5.9) -- node[above]{\texttt{app\_id}} (\xA,-5.9);
  \draw[ret] (\xD,-5.9) -- node[above]{\texttt{app\_id}} (\xE,-5.9);
\end{tikzpicture}
\caption[Application registration in an SD-QKD network]{Application registration in an SD-QKD network. Applications
know only their own key manager; the controller matches the two requests
and creates the service identifier. Adapted from~\cite[Fig.~2]{qkd015}.}
\label{fig:seq015}
\end{figure}


\paragraph{Physical link creation.}
The controller receives a trigger, from a demand detected on the nodes or
from its own NBI. It looks for two free, compatible interfaces at the two
sites, asks the optical network controller (SD-ONC) for an optical path
(\emph{lightpath}) with the constraints of the quantum channel and, once
confirmed, configures the two interfaces on their nodes. If the two
controllers are merged into one, the operation is internal. This workflow
is the core of the QKD–optical network coexistence demanded by UC4.

\paragraph{Virtual link creation.}
If there are no compatible interfaces or no suitable lightpath, the
controller falls back to a virtual link. It computes a path over the
existing QKD network, for example on a weighted graph built from key
availability, generation rates and consumption; it configures each
intermediate node in relay mode, indicating previous and next hop; finally
it configures the endpoints with the bandwidth to provide. The
specification does not prescribe the algorithm. A natural choice,
consistent with~\eqref{eq:eskr015}, is to weight each link with the inverse
of its available key and discard links that cannot sustain the requested
bandwidth $r$:
\begin{equation}
  w_\ell = \frac{1}{\mathrm{ESKR}_\ell}, \qquad
  p^\star = \arg\min_{p\,:\,\min_{\ell\in p}\mathrm{ESKR}_\ell \ge r}
            \;\sum_{\ell\in p} w_\ell .
  \label{eq:path015}
\end{equation}
Equation~\eqref{eq:path015} is our example, not part of the standard. The
specification warns that, if poorly managed, virtual links can congest the
underlying physical links: this is precisely why a controller with a global
view of consumption and availability is needed.

% ---------------------------------------------------------------------
\subsubsection{Protocols and security}

GS~QKD~015 does not mandate a protocol, but points to the two most common
YANG-based ones: NETCONF~\cite{rfc6241}, typically between device and
controller, transactional and carried over SSH or TLS with XML encoding;
and RESTCONF~\cite{rfc8040}, typically on the controllers' NBI, based on
HTTP with CRUD operations and JSON or XML encoding.

On security (clause 6) the specification distinguishes two cases when a new
system is installed. If the system joins a secure site already in the
network, the channel to the controller can be protected with QKD keys
produced by existing links. If a new secure site is added, an initial
scheme (a trusted identity or a pre-shared secret, used only during
installation) is needed to introduce the new SDN agent to the controller;
after the first authentication, QKD keys take over. It is the same
key-growing principle seen at the beginning.

% ---------------------------------------------------------------------
\subsubsection{From a single domain to orchestration}

GS~QKD~015 solves the control of \emph{one} QKD domain by \emph{one}
controller. Three problems already hinted at by GS~QKD~002 are left out.
First, an operator may have several QKD domains, possibly from different
vendors, each with its own controller. Second, the applications that
consume key may live in another network, for example in the nodes of the
optical transport network, managed by a different controller. Third, the
operator wants to manage QKD with the same orchestration tools it uses for
the rest of the network. A higher layer is therefore needed, together with
a standard interface to reach it: this is the job of GS~QKD~018.
